\documentclass[11pt]{article}
\usepackage[margin=25mm]{geometry}
\usepackage{fontspec}
\setmainfont{Latin Modern Roman}
\usepackage{amsmath,amssymb,mathtools}
\usepackage{unicode-math}
\setmathfont{Latin Modern Math}
\usepackage{xurl,hyperref}
\hypersetup{colorlinks=true,urlcolor=blue}
\setcounter{secnumdepth}{0}
\setlength{\parindent}{0pt}
\setlength{\parskip}{5pt}
\setlength{\emergencystretch}{3em}
\newcommand{\sref}[2]{\hyperlink{source:#1:#2}{[S#2]}}
\newcommand{\cataloguescope}[1]{\paragraph{Recorded target.}#1}
\newcommand{\R}{\mathbb R}
\newcommand{\C}{\mathbb C}
\newcommand{\N}{\mathbb N}
\newcommand{\Z}{\mathbb Z}
\newcommand{\Q}{\mathbb Q}
\newcommand{\F}{\mathbb F}
\newcommand{\eps}{\varepsilon}
\begin{document}
% UnsolvedMath ID 8100007; source catalogue code AMR-080-0007
\setcounter{section}{8100006}
\section{additional problems for integrable systems}\label{8100007}
{\small UnsolvedMath ID 8100007 · Mathematical Physics}
\subsection{Definitions and mathematical statement}

\paragraph{Shared notation.}$\mathbb N=\{1,2,\ldots\}$; the standard number fields have their usual meanings. All additional conventions are stated locally.


\paragraph{Definitions and notation.}A soliton is a localized persistent wave of an integrable PDE; a two-soliton initial profile produces two such waves in the unperturbed focusing NLS. In a Toda chain, particle accelerations are differences of adjacent exponential forces. A shock driver has $x_0(t)=2at$ and initially equally spaced particles at rest. For a trace-class positive contraction $K$, the Fredholm determinant $\det(I-K)=\prod_j(1-\lambda_j)$ is naturally interpretable through independent Bernoulli variables with probabilities $\lambda_j$. In the source's power perturbation, $\ell>0$ denotes the exponent and $\varepsilon>0$ is sufficiently small for the regime under consideration.

\paragraph{Catalogue wording retained for traceability.}
Analyze the long-time Cauchy problem for the perturbed focusing equation $i\widetilde u_t+\widetilde u_{xx}+2|\widetilde u|^2\widetilde u+\varepsilon|\widetilde u|^\ell\widetilde u=0$, with decaying initial data, especially data giving a two-soliton solution when $\varepsilon=0$. Prove rigorously by Riemann--Hilbert steepest descent the asymptotic binary periodic motion in the Toda shock problem with $x_k(0)=k$, $\dot x_k(0)=0$ for $k\ge1$, $x_0(t)=2at$, $a>1$. Analyze the same shock problem for suitable small perturbations $F$ of the force $e^x$, under the source’s condition that the corresponding doubly infinite chain has a temporally periodic, spatially two-periodic solution. Seek a useful intrinsic probabilistic definition of integrability, motivated by expressing key statistics through independent coin flips.
\paragraph{Scope and source qualification.}The printed focusing equation has an untilde u in its cubic term; the consistent unknown is explicitly used here as an editorial correction. The truncated integrable-probability question is restored. Suitability assumptions on perturbations are left as a required part of the programme, and the two-slit target has its separate canonical record 8100008. The source's defocusing theorem assumes $\ell>7/2$, while it expects extension to $\ell>2$; its later focusing comparison uses $\ell>2$ as the proposed regime. A proved defocusing threshold is not a proved focusing threshold.
\subsection{Short English statement}
Analyze perturbed focusing solitons and Toda shocks, and seek a definition of integrability intrinsic to probability.
\subsection{Research attempt}
\paragraph{Model, date and investigated scope.}
OpenAI GPT-6 Astra (configured model \texttt{gpt-6-astra}), ultra reasoning; generated 11 October 2026. Proved the independent-Bernoulli interpretation for every positive trace-class contraction and showed why that property alone is too broad to characterize integrability. All focusing and Toda asymptotic questions remain unresolved.

\paragraph{Sources and literature check.}
Read the supplied complete definition, catalogue wording, scope qualifications and source list for this canonical ID. The original-source citations are inherited as supplied references; no fresh primary-source or novelty search was performed for this attempt. No current open-status claim is inferred from this calculation.

\paragraph{Testing the independent-coin characterization.}Let \(K\) be any positive trace-class contraction, with eigenvalues \(0\le\lambda_j\le1\) and \(\sum_j\lambda_j<\infty\). Construct independent Bernoulli variables \(B_j\) with \(\mathbb P(B_j=1)=\lambda_j\). Then \(S=\sum_jB_j\) is finite almost surely: its expectation is the finite sum of the probabilities, and a nonnegative random variable with finite expectation cannot be infinite on a positive-probability event. Independence and finite products give
\[\mathbb E z^S=\prod_j(1-\lambda_j+z\lambda_j),\quad0\le z\le1,\qquad
\mathbb P(S=0)=\prod_j(1-\lambda_j)=\det(I-K).
\]
The infinite-product formula follows by bounded convergence from partial sums. This proves the desired coin interpretation at the spectral level for an arbitrary positive trace-class contraction.
\paragraph{A limitation of this definition.}The property therefore does not by itself distinguish an algebraically solvable model: the eigenvalues of \(K\) may be inaccessible, and only a single count distribution is represented. Conversely any prescribed summable sequence in \([0,1]\) defines a diagonal positive trace-class contraction and hence such a coin model. A useful intrinsic definition needs constraints on how the statistic and its kernel are generated, and on accessible joint observables, not merely the existence of this product after diagonalization.
\paragraph{Full target remaining.}No long-time theorem is proved for the perturbed focusing two-soliton equation, no binary-periodic Toda shock asymptotic is derived, and no perturbation of the Toda force is controlled under the stated two-periodic-solution hypothesis. The source's focusing exponent regime is kept distinct from its proved defocusing context. This elementary probability calculation addresses one proposed criterion and exposes its insufficiency; it is not a definition resolving the complete integrability programme.

\paragraph{Verification and adversarial audit.}
Checked trace summability, almost-sure finiteness, bounded-convergence product limit, eigenvalue-one case and all bundled focusing/Toda scopes. This was an LLM adversarial audit, not independent human verification; no theorem prover was used. Prompts A then B applied to this ID: Checked trace summability, almost-sure finiteness, bounded-convergence product limit, eigenvalue-one case and all bundled focusing/Toda scopes. Scope remains limited: Add an intrinsic, non-tautological structural criterion for probabilistic integrability; separately establish the required nonlinear long-time estimates.

Coordinator checked almost-sure finiteness of the Bernoulli sum and bounded-convergence justification of the spectral product. This additional check was model-assisted, not independent human verification.

\paragraph{Tools, process and contributions.}
Python 3 and local shell used to read supplied files and serialize this manuscript; mathematics derived symbolically in the text, with no numerical experiment or formal proof assistant. Contribution prompt A was followed by prompt B. GPT-6 Astra with ultra reasoning drafted the research and performed its stated self-audit; any additional coordinator contribution is identified in the verification and audit. A GPT-6-based Codex coordinator assembled the manuscript and checked repository integration; its exact variant and reasoning configuration were not exposed. No human mathematical contribution, independent human proof check or formal proof certification is supplied. Complete derivations and the consulted source references are included here; private reasoning transcripts are not published.

\paragraph{Final status and first remaining gap.}
\textbf{UNRESOLVED}. Add an intrinsic, non-tautological structural criterion for probabilistic integrability; separately establish the required nonlinear long-time estimates.

\section{Completion Estimate}
\noindent\textbf{COMPLETION: 0\%}
The calculation is an elementary benchmark or conditional reduction; it establishes no claimed new portion of the complete source target.

\subsection{Sources}
\begin{itemize}
\item[\hypertarget{source:8100007:1}{S1}] ULAM.AI and the UnsolvedMath Contributors, UnsolvedMath, record 8100007 (AMR-080-0007). Catalogue curation under CC BY 4.0; definitions and English exposition added locally.
\url{https://huggingface.co/datasets/ulamai/UnsolvedMath}.
\item[\hypertarget{source:8100007:2}{S2}] Percy Deift, Some Open Problems in Random Matrix Theory and the Theory of Integrable Systems. II, SIGMA 13 (2017), 016, Problem 7.
\url{https://arxiv.org/abs/1703.04931}.
\item[\hypertarget{source:8100007:3}{S3}] Primary full-text formulation and 2017 historical progress discussion.
\url{https://arxiv.org/html/1703.04931v1}.
\end{itemize}
\label{end:8100007}
\end{document}
